\section{Security and Trust Boundaries}
\label{sec:trust}

The platform has several layers, each with distinct authentication and trust assumptions. This section defines their guarantees, boundaries, and excluded threats.

\subsection{Transport: Hop-by-Hop Encryption and Authentication}

A connection between two neighboring peers is established by an authenticated ephemeral Diffie-Hellman handshake (X25519). The handshake exchanges symmetric keys used to encrypt subsequent traffic with AEAD constructions: AES-128-GCM by default, with ChaCha20-Poly1305 as an alternative for platforms without AES hardware acceleration. The handshake itself is encrypted with the recipient's long-term identity key, which doubles as a defense against unsolicited probes: a peer that has not received a properly-encrypted hello does not respond at all, making the port effectively invisible to a passive scanner.

This authentication is \emph{hop-by-hop}. A request that traverses peers $A \to B \to C \to D$ is authenticated by $A$ to $B$, by $B$ to $C$, and by $C$ to $D$. There is no transport-level guarantee that the message body at $D$ is the body $A$ originally sent: an intermediate peer can drop, reorder, or modify the payload before forwarding.

\subsection{Application Layer: End-to-End Authentication via Validity Predicates}

End-to-end integrity is the contract's responsibility. A contract's validity predicate $\valid_\contract$ runs at every peer that touches the state. Update messages typically include a signature by an application-authorized key over the update content; the validity predicate verifies the signature and rejects the update otherwise. Because $\valid_\contract$ is itself part of the contract's content-addressed code, a peer cannot serve a state that fails validity without being caught by every receiving peer that runs the same code.

The combination is the platform's substitute for a single end-to-end transport channel. Transport gives hop-by-hop confidentiality and authentication of neighbors. The contract gives end-to-end semantic integrity of state.

\subsection{NAT Traversal and Bootstrap}

Most peers operate behind consumer NAT and do not have a stable inbound address. The transport supports hole-punching: when two peers want to connect and at least one is behind a NAT, both peers send hello packets toward each other's observed external addresses simultaneously \citep{ford2005peer}. The introducer in this scheme, the peer that reports each side's externally-observed address, is whichever peer routed the connection request between them.

For a brand-new peer joining the network the introducer is a gateway, which is required to have a public IP. For a peer that already has at least one connection, the introducer is one of its existing neighbors. Gateways are required only for the initial join. Existing peers introduce all subsequent connections, so the network can continue forming connections while every gateway is offline.

A peer's location is derived from this externally-observed address. The implications are discussed in the next subsection.

\subsection{Location Assignment: A Known Sybil and Grinding Surface}

A peer's location is computed by hashing its externally-observed network address into $[0, 1)$, but the hash is taken over the address \emph{prefix} rather than the full address: the high 24 bits of an IPv4 address (so all hosts in the same /24 share a location) and the high 48 bits of an IPv6 address (the typical site-level allocation). The intent is to raise the cost of location grinding: an adversary that wants to place peers at chosen ring locations needs control of multiple distinct prefixes, not merely many hosts inside one subnet.

Prefix hashing only raises the cost of the attack. An adversary who can rent IP space across many /24s, obtain an IPv6 range, or provision VMs across distinct cloud prefixes can still grind locations and concentrate peers in chosen ring regions. Those peers can selectively answer contract requests or crowd honest peers out under the topology rule.

Prefix hashing also places every honest peer sharing a prefix at the \emph{same} ring location. Independently operated peers behind a university or corporate NAT, a carrier-grade NAT pool, or a popular cloud region consequently have ring distance zero from one another. This common deployment case affects topology construction, replication, and the application-layer Sybil mitigations described below.

\emph{Topology construction.} Gap-targeting (\cref{sec:routing}) assumes that neighbor locations are distributed around the ring. A cluster of co-located peers appears as a zero-width spike to the largest-log-distance-gap computation. \emph{Replication.} Peers near a contract location are intended to protect against a single peer serving stale or doctored state. Co-located peers may instead share an operator and failure domain, reducing the independence of those replicas. \emph{Sybil bootstrapping.} Application-layer defenses depend on honest peers being geographically diverse near critical locations, which shared prefixes make less likely.

Adding a per-peer nonce would spread co-located peers but would also restore the grinding freedom that prefix hashing is intended to reduce. Resolving this tension remains open (\cref{sec:status}).

The platform does not solve this at the location layer. Two architectural mitigations operate at higher levels.

First, the contract layer's validity predicates make state-doctoring detectable: a malicious peer that returns a state with an invalid signature is rejected by every honest receiver. The attack surface this leaves open is selectively serving stale (but still valid) state, not actively forging state.

Second, the platform expects identity, reputation, and contribution-control to be built as application-layer contracts (see Ghost Keys \citep{ghostkey} for an example of a contract-mediated anonymous identity scheme). A contract that posts authoritative information can require its writers to hold reputation tokens issued by another contract; an inbox contract can demand a small proof of work; a community forum can require introductions from existing members. The platform's role is to provide a substrate on which these mechanisms can be built, not to embed any particular one in the routing layer.

The current design leaves Sybil defense to applications. A platform-level scheme would either impose a particular economic or identity model or provide only a generic obstacle that a sufficiently resourced adversary could overcome.

Application-layer contracts are themselves routed and replicated through the vulnerable location layer. An attacker with enough /24 prefixes to crowd a contract's location can serve poisoned reads and prevent honest writes from reaching its canonical replicas. Validity predicates catch forgery but not selective service of stale, valid state. Bootstrapping therefore requires a property the platform does not currently provide: geographic diversity near safety-critical contracts, sufficient ring-distance separation among their replicas, or out-of-band trust anchors such as a bundled root key or release-managed seed list. Selecting and implementing such a mechanism remains open work.

\subsection{Delegates: A Trust Boundary Around Secrets}

Contracts hold public state. Applications also have private state (a user's signing keys, a private contact list, the decryption material for an end-to-end encrypted inbox) that should never be replicated and should be accessible only through controlled interfaces. The platform's primitive for this is the delegate.

A delegate is a WebAssembly module running inside the local Freenet Core on a user's device. It holds private secret state, sandboxed by the core. Other components (user interfaces, contracts, other delegates) can interact with it only by sending messages. The core attaches a verified sender identifier to each inbound message, so the delegate's policy can be ``only respond to messages from UI X, originating from contract Y.''

The trust boundary is therefore platform-enforced rather than conventional. Traditional in-process encapsulation (a ``private'' field on an object) is honored only by code that chooses to honor it; the same address space affords access to memory, file system, and undocumented APIs. A delegate's secret state is not addressable from any other code on the device. The only way to extract a secret from a delegate is to convince the delegate, on the basis of sender identity and message contents, that releasing the secret is allowed.

The boundary is enforced by the local Freenet Core, which is itself a piece of software running on the same host. A compromised or tampered Core can violate the boundary directly. Defending against this requires hardware-attested execution or TEE-style mechanisms that the platform does not currently provide; the trust model here assumes an uncompromised local Core.

\paragraph{Worked example: signing without exporting a key.} A chat application's UI needs to send signed messages. Without delegates, it must hold the private key. With delegates, the UI sends a $(\textsf{sign}, \mathrm{payload})$ message to a key-manager delegate; the delegate verifies the sender, signs the payload, and returns the signature. The private key never crosses the boundary. A bug or compromise in the UI can produce unwanted signed payloads but cannot exfiltrate the key itself; the delegate is also free to apply rate limits or user confirmation prompts before signing.

\paragraph{Cross-device replication.} A user typically runs the ``same'' delegate on multiple devices, and the platform's intended pattern for keeping those instances in sync is a shared-secret contract: a contract that holds an encrypted, replicated representation of the shared private state, decryptable only by the delegates that hold the symmetric key. The contract's merge function provides convergence; the delegate provides confidentiality. This pattern is only partially implemented at the time of writing (see \cref{sec:status}); the encrypted-contract primitives exist but the end-to-end delegate-to-delegate flow is not yet wired up.

\subsection{Resource Accounting and What It Doesn't Cover}

Each peer enforces limits on bandwidth (token-bucket rate limiting at the transport layer), storage (a byte-budgeted set of hosted contracts), CPU, and memory (per-contract WebAssembly fuel and explicit memory bounds, enforced by Wasmtime \citep{wasmtime}). The storage budget is denominated in hosted-state bytes and defaults to a clamped fraction of host memory because resident memory and update fan-out constrain the number of active contracts more than disk space does. An operator may override the default. When over budget, a peer removes the contract with the fewest dependent subscribers, breaking ties by least-recent genuine access. Excessive resource consumption by a neighbor results in disconnection.

These limits protect against a single misbehaving neighbor, but not coordinated traffic from many neighbors or denial of service through many cheap identities. Addressing those attacks probably requires accounting that persists across identities, as discussed in \cref{sec:status}.
